% --- Archon physics formalization source begin ---
% archon:physics
% archon:covers PhyXMiniProblems/problem_phyx_mini_0228.lean
% archon:source-report reports/phyx_mini/problem_phyx_mini_0228.source.json

\chapter{Physics problem phyx\_mini\_0228}
\label{ch:PhyXMiniProblems_problem_phyx_mini_0228}

\paragraph{Problem source.}
\#\# Physical scenario

A watch balance wheel (figure) has a period of oscillation of \textbackslash{}( 0.250\textbackslash{},\textbackslash{}mathrm\{s\} \textbackslash{}). The wheel is constructed so that its mass of \textbackslash{}( 20.0\textbackslash{},\textbackslash{}mathrm\{g\} \textbackslash{}) is concentrated around a rim of radius \textbackslash{}( 0.500\textbackslash{},\textbackslash{}mathrm\{cm\} \textbackslash{}).

\#\# Question

What is the torsion constant of the attached spring?

\#\# Answer choices

- A: A: \$3.61\textbackslash{}times10\textasciicircum{}\{-4\}\textbackslash{}

- B: B: \$3.16\textbackslash{}times10\textasciicircum{}\{-3\}\textbackslash{}

- C: C: \$2.16\textbackslash{}times10\textasciicircum{}\{-4\}\textbackslash{}

- D: D: \$3.16\textbackslash{}times10\textasciicircum{}\{-4\}\textbackslash{}

\#\# Figure caption (auxiliary; use the image as primary evidence)

The image shows the inner mechanism of a watch, highlighting several intricate components. The main feature pointed out is the "Balance wheel," which is a critical part of watches that helps regulate the timekeeping. The mechanism includes multiple interconnected gears, wheels, and levers arranged in a circular fashion. The parts appear to be metallic and finely detailed, indicating precision engineering. Above the image, a label with a line points directly to the balance wheel, marking its location within the watch mechanism.

\#\# Dataset metadata

Resonance and Harmonics; Physical Model Grounding Reasoning, Multi-Formula Reasoning

\paragraph{Recorded answer/context.}
D: D: \$3.16\textbackslash{}times10\textasciicircum{}\{-4\}\textbackslash{}

\paragraph{Figure/image path.}
/root/proposal\_for\_physic/science-mango/phyx\_mini\_run/phyx\_data/test\_image/228.png

\paragraph{Formalization target.}
create a compiling Lean file with sorry bodies at `PhyXMiniProblems/problem\_phyx\_mini\_0228.lean`. The Lean declarations must preserve the physical quantities, dimensions or dimensional roles, figure labels, governing-law hypotheses, and final relation expressed by this problem.
Use Mathlib/Physlib names found through LeanExplore where available. If a physics API is missing, introduce faithful local abstractions rather than scalar placeholder aliases.

\begin{theorem}[Physics formalization target]
\label{thm:physics:phyx_mini_0228:target}
\lean{PhyXMiniProblems.ProblemPhyXMini0228.problem_phyx_mini_0228}
\uses{def:physics:phyx-mini-0228:phyxminiproblems-problemphyxmini0228-torsionconstantinsi, def:physics:phyx-mini-0228:phyxminiproblems-problemphyxmini0228-watchbalancewheelsetup, def:physics:phyx-mini-0228:phyxminiproblems-problemphyxmini0228-matchesproblemandfigurereadouts, def:physics:phyx-mini-0228:phyxminiproblems-problemphyxmini0228-usesidealtorsionaloscillatormodel, def:physics:phyx-mini-0228:phyxminiproblems-problemphyxmini0228-hasphysicalbalancewheelparameters, def:physics:phyx-mini-0228:phyxminiproblems-problemphyxmini0228-satisfiesrimmomentofinertialaw, def:physics:phyx-mini-0228:phyxminiproblems-problemphyxmini0228-satisfiestorsionaloscillatorperiodlaw, def:physics:phyx-mini-0228:phyxminiproblems-problemphyxmini0228-recordeddatasetanswer, def:physics:phyx-mini-0228:phyxminiproblems-problemphyxmini0228-matchesanswerchoice}
The assigned autoformalize agent should translate this physics problem into Lean declarations in the covered file, with theorem and lemma proof bodies written as `by sorry`.
\end{theorem}
\begin{proof}
This is an autoformalization task, not a proof task. Produce statements that can later be proved by the physics prover without weakening the physical model.
\end{proof}

% --- Archon declaration topology begin ---
\section{Declaration topology}
\begin{definition}[Mass Quantity]
  \label{def:physics:phyx-mini-0228:phyxminiproblems-problemphyxmini0228-massquantity}
  \lean{PhyXMiniProblems.ProblemPhyXMini0228.MassQuantity}
  A nonnegative physical mass, independent of the unit used to read it
\end{definition}
\begin{proof}
  This is the declared model definition of Mass Quantity.
\end{proof}
\begin{definition}[Length Quantity]
  \label{def:physics:phyx-mini-0228:phyxminiproblems-problemphyxmini0228-lengthquantity}
  \lean{PhyXMiniProblems.ProblemPhyXMini0228.LengthQuantity}
  A nonnegative physical length, used for the balance-wheel rim radius
\end{definition}
\begin{proof}
  This is the declared model definition of Length Quantity.
\end{proof}
\begin{definition}[Time Quantity]
  \label{def:physics:phyx-mini-0228:phyxminiproblems-problemphyxmini0228-timequantity}
  \lean{PhyXMiniProblems.ProblemPhyXMini0228.TimeQuantity}
  A nonnegative physical duration, used for the oscillation period
\end{definition}
\begin{proof}
  This is the declared model definition of Time Quantity.
\end{proof}
\begin{definition}[Moment Of Inertia Quantity]
  \label{def:physics:phyx-mini-0228:phyxminiproblems-problemphyxmini0228-momentofinertiaquantity}
  \lean{PhyXMiniProblems.ProblemPhyXMini0228.MomentOfInertiaQuantity}
  Moment of inertia, carrying the physical dimension mass times length squared
\end{definition}
\begin{proof}
  This is the declared model definition of Moment Of Inertia Quantity.
\end{proof}
\begin{definition}[Torsion Constant Quantity]
  \label{def:physics:phyx-mini-0228:phyxminiproblems-problemphyxmini0228-torsionconstantquantity}
  \lean{PhyXMiniProblems.ProblemPhyXMini0228.TorsionConstantQuantity}
  Torsion constant of a rotational spring. Radians are dimensionless, so N m / rad has dimension mass * length\textasciicircum{}2 / time\textasciicircum{}2, the same dimension as energy but a distinct physical role in this model
\end{definition}
\begin{proof}
  This is the declared model definition of Torsion Constant Quantity.
\end{proof}
\begin{definition}[mass Readout]
  \label{def:physics:phyx-mini-0228:phyxminiproblems-problemphyxmini0228-massreadout}
  \lean{PhyXMiniProblems.ProblemPhyXMini0228.massReadout}
  \uses{def:physics:phyx-mini-0228:phyxminiproblems-problemphyxmini0228-massquantity}
  Read a physical mass in a selected mass unit
\end{definition}
\begin{proof}
  This is the declared model definition of mass Readout.
\end{proof}
\begin{definition}[length Readout]
  \label{def:physics:phyx-mini-0228:phyxminiproblems-problemphyxmini0228-lengthreadout}
  \lean{PhyXMiniProblems.ProblemPhyXMini0228.lengthReadout}
  \uses{def:physics:phyx-mini-0228:phyxminiproblems-problemphyxmini0228-lengthquantity}
  Read a physical length in a selected length unit
\end{definition}
\begin{proof}
  This is the declared model definition of length Readout.
\end{proof}
\begin{definition}[time Readout]
  \label{def:physics:phyx-mini-0228:phyxminiproblems-problemphyxmini0228-timereadout}
  \lean{PhyXMiniProblems.ProblemPhyXMini0228.timeReadout}
  \uses{def:physics:phyx-mini-0228:phyxminiproblems-problemphyxmini0228-timequantity}
  Read a physical duration in a selected time unit
\end{definition}
\begin{proof}
  This is the declared model definition of time Readout.
\end{proof}
\begin{definition}[moment Of Inertia Readout]
  \label{def:physics:phyx-mini-0228:phyxminiproblems-problemphyxmini0228-momentofinertiareadout}
  \lean{PhyXMiniProblems.ProblemPhyXMini0228.momentOfInertiaReadout}
  \uses{def:physics:phyx-mini-0228:phyxminiproblems-problemphyxmini0228-momentofinertiaquantity}
  Read moment of inertia in the selected mass unit times length unit squared
\end{definition}
\begin{proof}
  This is the declared model definition of moment Of Inertia Readout.
\end{proof}
\begin{definition}[torsion Constant Readout]
  \label{def:physics:phyx-mini-0228:phyxminiproblems-problemphyxmini0228-torsionconstantreadout}
  \lean{PhyXMiniProblems.ProblemPhyXMini0228.torsionConstantReadout}
  \uses{def:physics:phyx-mini-0228:phyxminiproblems-problemphyxmini0228-torsionconstantquantity}
  Read a torsion constant in the coherent unit massUnit * lengthUnit\textasciicircum{}2 / timeUnit\textasciicircum{}2 per dimensionless radian
\end{definition}
\begin{proof}
  This is the declared model definition of torsion Constant Readout.
\end{proof}
\begin{definition}[period In Seconds]
  \label{def:physics:phyx-mini-0228:phyxminiproblems-problemphyxmini0228-periodinseconds}
  \lean{PhyXMiniProblems.ProblemPhyXMini0228.periodInSeconds}
  \uses{def:physics:phyx-mini-0228:phyxminiproblems-problemphyxmini0228-timequantity, def:physics:phyx-mini-0228:phyxminiproblems-problemphyxmini0228-timereadout}
  The period readout in seconds
\end{definition}
\begin{proof}
  This is the declared model definition of period In Seconds.
\end{proof}
\begin{definition}[mass In Grams]
  \label{def:physics:phyx-mini-0228:phyxminiproblems-problemphyxmini0228-massingrams}
  \lean{PhyXMiniProblems.ProblemPhyXMini0228.massInGrams}
  \uses{def:physics:phyx-mini-0228:phyxminiproblems-problemphyxmini0228-massquantity, def:physics:phyx-mini-0228:phyxminiproblems-problemphyxmini0228-massreadout}
  The wheel-mass readout in grams
\end{definition}
\begin{proof}
  This is the declared model definition of mass In Grams.
\end{proof}
\begin{definition}[length In Centimeters]
  \label{def:physics:phyx-mini-0228:phyxminiproblems-problemphyxmini0228-lengthincentimeters}
  \lean{PhyXMiniProblems.ProblemPhyXMini0228.lengthInCentimeters}
  \uses{def:physics:phyx-mini-0228:phyxminiproblems-problemphyxmini0228-lengthquantity, def:physics:phyx-mini-0228:phyxminiproblems-problemphyxmini0228-lengthreadout}
  The rim-radius readout in centimeters
\end{definition}
\begin{proof}
  This is the declared model definition of length In Centimeters.
\end{proof}
\begin{definition}[mass In Kilograms]
  \label{def:physics:phyx-mini-0228:phyxminiproblems-problemphyxmini0228-massinkilograms}
  \lean{PhyXMiniProblems.ProblemPhyXMini0228.massInKilograms}
  \uses{def:physics:phyx-mini-0228:phyxminiproblems-problemphyxmini0228-massquantity, def:physics:phyx-mini-0228:phyxminiproblems-problemphyxmini0228-massreadout}
  The wheel-mass readout in SI kilograms
\end{definition}
\begin{proof}
  This is the declared model definition of mass In Kilograms.
\end{proof}
\begin{definition}[length In Meters]
  \label{def:physics:phyx-mini-0228:phyxminiproblems-problemphyxmini0228-lengthinmeters}
  \lean{PhyXMiniProblems.ProblemPhyXMini0228.lengthInMeters}
  \uses{def:physics:phyx-mini-0228:phyxminiproblems-problemphyxmini0228-lengthquantity, def:physics:phyx-mini-0228:phyxminiproblems-problemphyxmini0228-lengthreadout}
  The rim-radius readout in SI meters
\end{definition}
\begin{proof}
  This is the declared model definition of length In Meters.
\end{proof}
\begin{definition}[moment Of Inertia In SI]
  \label{def:physics:phyx-mini-0228:phyxminiproblems-problemphyxmini0228-momentofinertiainsi}
  \lean{PhyXMiniProblems.ProblemPhyXMini0228.momentOfInertiaInSI}
  \uses{def:physics:phyx-mini-0228:phyxminiproblems-problemphyxmini0228-momentofinertiaquantity, def:physics:phyx-mini-0228:phyxminiproblems-problemphyxmini0228-momentofinertiareadout}
  Moment of inertia in SI kg m\textasciicircum{}2
\end{definition}
\begin{proof}
  This is the declared model definition of moment Of Inertia In SI.
\end{proof}
\begin{definition}[torsion Constant In SI]
  \label{def:physics:phyx-mini-0228:phyxminiproblems-problemphyxmini0228-torsionconstantinsi}
  \lean{PhyXMiniProblems.ProblemPhyXMini0228.torsionConstantInSI}
  \uses{def:physics:phyx-mini-0228:phyxminiproblems-problemphyxmini0228-torsionconstantquantity, def:physics:phyx-mini-0228:phyxminiproblems-problemphyxmini0228-torsionconstantreadout}
  Torsion constant in SI N m / rad, equivalently kg m\textasciicircum{}2 s\textasciicircum{}-2
\end{definition}
\begin{proof}
  This is the declared model definition of torsion Constant In SI.
\end{proof}
\begin{definition}[Watch Component]
  \label{def:physics:phyx-mini-0228:phyxminiproblems-problemphyxmini0228-watchcomponent}
  \lean{PhyXMiniProblems.ProblemPhyXMini0228.WatchComponent}
  Watch components distinguished by the source and the primary image
\end{definition}
\begin{proof}
  This is the declared model definition of Watch Component.
\end{proof}
\begin{definition}[Figure Label]
  \label{def:physics:phyx-mini-0228:phyxminiproblems-problemphyxmini0228-figurelabel}
  \lean{PhyXMiniProblems.ProblemPhyXMini0228.FigureLabel}
  Text labels visible in the supplied image
\end{definition}
\begin{proof}
  This is the declared model definition of Figure Label.
\end{proof}
\begin{definition}[Wheel Mass Distribution]
  \label{def:physics:phyx-mini-0228:phyxminiproblems-problemphyxmini0228-wheelmassdistribution}
  \lean{PhyXMiniProblems.ProblemPhyXMini0228.WheelMassDistribution}
  The radial mass-distribution idealization used for the wheel
\end{definition}
\begin{proof}
  This is the declared model definition of Wheel Mass Distribution.
\end{proof}
\begin{definition}[Attached Spring Kind]
  \label{def:physics:phyx-mini-0228:phyxminiproblems-problemphyxmini0228-attachedspringkind}
  \lean{PhyXMiniProblems.ProblemPhyXMini0228.AttachedSpringKind}
  The mechanical type of the spring attached to the balance wheel
\end{definition}
\begin{proof}
  This is the declared model definition of Attached Spring Kind.
\end{proof}
\begin{definition}[Oscillation Regime]
  \label{def:physics:phyx-mini-0228:phyxminiproblems-problemphyxmini0228-oscillationregime}
  \lean{PhyXMiniProblems.ProblemPhyXMini0228.OscillationRegime}
  The ideal oscillation regime in which the standard torsional period law applies
\end{definition}
\begin{proof}
  This is the declared model definition of Oscillation Regime.
\end{proof}
\begin{definition}[Balance Wheel Figure]
  \label{def:physics:phyx-mini-0228:phyxminiproblems-problemphyxmini0228-balancewheelfigure}
  \lean{PhyXMiniProblems.ProblemPhyXMini0228.BalanceWheelFigure}
  \uses{def:physics:phyx-mini-0228:phyxminiproblems-problemphyxmini0228-watchcomponent, def:physics:phyx-mini-0228:phyxminiproblems-problemphyxmini0228-figurelabel}
  Qualitative evidence read from the supplied watch image
\end{definition}
\begin{proof}
  This is the declared model definition of Balance Wheel Figure.
\end{proof}
\begin{definition}[Watch Balance Wheel Setup]
  \label{def:physics:phyx-mini-0228:phyxminiproblems-problemphyxmini0228-watchbalancewheelsetup}
  \lean{PhyXMiniProblems.ProblemPhyXMini0228.WatchBalanceWheelSetup}
  \uses{def:physics:phyx-mini-0228:phyxminiproblems-problemphyxmini0228-massquantity, def:physics:phyx-mini-0228:phyxminiproblems-problemphyxmini0228-lengthquantity, def:physics:phyx-mini-0228:phyxminiproblems-problemphyxmini0228-timequantity, def:physics:phyx-mini-0228:phyxminiproblems-problemphyxmini0228-momentofinertiaquantity, def:physics:phyx-mini-0228:phyxminiproblems-problemphyxmini0228-torsionconstantquantity, def:physics:phyx-mini-0228:phyxminiproblems-problemphyxmini0228-wheelmassdistribution, def:physics:phyx-mini-0228:phyxminiproblems-problemphyxmini0228-attachedspringkind, def:physics:phyx-mini-0228:phyxminiproblems-problemphyxmini0228-oscillationregime, def:physics:phyx-mini-0228:phyxminiproblems-problemphyxmini0228-balancewheelfigure}
  The physical quantities are independent fields. In particular, springTorsionConstant is not defined from the answer value or from the other measurements; it is constrained only by the governing torsional-oscillator law
\end{definition}
\begin{proof}
  This is the declared model definition of Watch Balance Wheel Setup.
\end{proof}
\begin{definition}[Matches Problem And Figure Readouts]
  \label{def:physics:phyx-mini-0228:phyxminiproblems-problemphyxmini0228-matchesproblemandfigurereadouts}
  \lean{PhyXMiniProblems.ProblemPhyXMini0228.MatchesProblemAndFigureReadouts}
  \uses{def:physics:phyx-mini-0228:phyxminiproblems-problemphyxmini0228-periodinseconds, def:physics:phyx-mini-0228:phyxminiproblems-problemphyxmini0228-massingrams, def:physics:phyx-mini-0228:phyxminiproblems-problemphyxmini0228-lengthincentimeters, def:physics:phyx-mini-0228:phyxminiproblems-problemphyxmini0228-watchbalancewheelsetup}
  Numerical data stated in the problem and qualitative evidence visible in the primary image. This predicate contains no torsion-constant value and names no answer choice
\end{definition}
\begin{proof}
  This is the declared model definition of Matches Problem And Figure Readouts.
\end{proof}
\begin{definition}[Uses Ideal Torsional Oscillator Model]
  \label{def:physics:phyx-mini-0228:phyxminiproblems-problemphyxmini0228-usesidealtorsionaloscillatormodel}
  \lean{PhyXMiniProblems.ProblemPhyXMini0228.UsesIdealTorsionalOscillatorModel}
  \uses{def:physics:phyx-mini-0228:phyxminiproblems-problemphyxmini0228-watchbalancewheelsetup}
  The idealized small, undamped torsional-oscillation regime
\end{definition}
\begin{proof}
  This is the declared model definition of Uses Ideal Torsional Oscillator Model.
\end{proof}
\begin{definition}[Has Physical Balance Wheel Parameters]
  \label{def:physics:phyx-mini-0228:phyxminiproblems-problemphyxmini0228-hasphysicalbalancewheelparameters}
  \lean{PhyXMiniProblems.ProblemPhyXMini0228.HasPhysicalBalanceWheelParameters}
  \uses{def:physics:phyx-mini-0228:phyxminiproblems-problemphyxmini0228-massreadout, def:physics:phyx-mini-0228:phyxminiproblems-problemphyxmini0228-lengthreadout, def:physics:phyx-mini-0228:phyxminiproblems-problemphyxmini0228-timereadout, def:physics:phyx-mini-0228:phyxminiproblems-problemphyxmini0228-momentofinertiareadout, def:physics:phyx-mini-0228:phyxminiproblems-problemphyxmini0228-torsionconstantreadout, def:physics:phyx-mini-0228:phyxminiproblems-problemphyxmini0228-watchbalancewheelsetup}
  Positivity and nondegeneracy conditions for the mechanical parameters
\end{definition}
\begin{proof}
  This is the declared model definition of Has Physical Balance Wheel Parameters.
\end{proof}
\begin{definition}[Satisfies Rim Moment Of Inertia Law]
  \label{def:physics:phyx-mini-0228:phyxminiproblems-problemphyxmini0228-satisfiesrimmomentofinertialaw}
  \lean{PhyXMiniProblems.ProblemPhyXMini0228.SatisfiesRimMomentOfInertiaLaw}
  \uses{def:physics:phyx-mini-0228:phyxminiproblems-problemphyxmini0228-massreadout, def:physics:phyx-mini-0228:phyxminiproblems-problemphyxmini0228-lengthreadout, def:physics:phyx-mini-0228:phyxminiproblems-problemphyxmini0228-momentofinertiareadout, def:physics:phyx-mini-0228:phyxminiproblems-problemphyxmini0228-watchbalancewheelsetup}
  Thin-ring moment-of-inertia law for a mass concentrated around a rim: I = m r\textasciicircum{}2, stated in every coherent pair of mass and length units
\end{definition}
\begin{proof}
  This is the declared model definition of Satisfies Rim Moment Of Inertia Law.
\end{proof}
\begin{definition}[Satisfies Torsional Oscillator Period Law]
  \label{def:physics:phyx-mini-0228:phyxminiproblems-problemphyxmini0228-satisfiestorsionaloscillatorperiodlaw}
  \lean{PhyXMiniProblems.ProblemPhyXMini0228.SatisfiesTorsionalOscillatorPeriodLaw}
  \uses{def:physics:phyx-mini-0228:phyxminiproblems-problemphyxmini0228-timereadout, def:physics:phyx-mini-0228:phyxminiproblems-problemphyxmini0228-momentofinertiareadout, def:physics:phyx-mini-0228:phyxminiproblems-problemphyxmini0228-torsionconstantreadout, def:physics:phyx-mini-0228:phyxminiproblems-problemphyxmini0228-watchbalancewheelsetup}
  Governing period law for an ideal torsional oscillator: T = 2 pi sqrt (I / kappa). The law relates the unknown torsion constant to the independently stored period and inertia, but assigns it no answer value
\end{definition}
\begin{proof}
  This is the declared model definition of Satisfies Torsional Oscillator Period Law.
\end{proof}
\begin{definition}[Answer Choice]
  \label{def:physics:phyx-mini-0228:phyxminiproblems-problemphyxmini0228-answerchoice}
  \lean{PhyXMiniProblems.ProblemPhyXMini0228.AnswerChoice}
  Labels of the four numerical choices printed with the problem
\end{definition}
\begin{proof}
  This is the declared model definition of Answer Choice.
\end{proof}
\begin{definition}[displayed Answer Torsion Constant SI]
  \label{def:physics:phyx-mini-0228:phyxminiproblems-problemphyxmini0228-displayedanswertorsionconstantsi}
  \lean{PhyXMiniProblems.ProblemPhyXMini0228.displayedAnswerTorsionConstantSI}
  \uses{def:physics:phyx-mini-0228:phyxminiproblems-problemphyxmini0228-answerchoice}
  Torsion-constant readout printed beside each choice, interpreted in N m / rad
\end{definition}
\begin{proof}
  This is the declared model definition of displayed Answer Torsion Constant SI.
\end{proof}
\begin{definition}[displayed Answer Tolerance SI]
  \label{def:physics:phyx-mini-0228:phyxminiproblems-problemphyxmini0228-displayedanswertolerancesi}
  \lean{PhyXMiniProblems.ProblemPhyXMini0228.displayedAnswerToleranceSI}
  \uses{def:physics:phyx-mini-0228:phyxminiproblems-problemphyxmini0228-answerchoice}
  Half a unit in the final displayed significant digit of each choice
\end{definition}
\begin{proof}
  This is the declared model definition of displayed Answer Tolerance SI.
\end{proof}
\begin{definition}[recorded Dataset Answer]
  \label{def:physics:phyx-mini-0228:phyxminiproblems-problemphyxmini0228-recordeddatasetanswer}
  \lean{PhyXMiniProblems.ProblemPhyXMini0228.recordedDatasetAnswer}
  \uses{def:physics:phyx-mini-0228:phyxminiproblems-problemphyxmini0228-answerchoice}
  The answer label recorded by the supplied dataset metadata
\end{definition}
\begin{proof}
  This is the declared model definition of recorded Dataset Answer.
\end{proof}
\begin{definition}[Matches Answer Choice]
  \label{def:physics:phyx-mini-0228:phyxminiproblems-problemphyxmini0228-matchesanswerchoice}
  \lean{PhyXMiniProblems.ProblemPhyXMini0228.MatchesAnswerChoice}
  \uses{def:physics:phyx-mini-0228:phyxminiproblems-problemphyxmini0228-torsionconstantquantity, def:physics:phyx-mini-0228:phyxminiproblems-problemphyxmini0228-torsionconstantinsi, def:physics:phyx-mini-0228:phyxminiproblems-problemphyxmini0228-answerchoice, def:physics:phyx-mini-0228:phyxminiproblems-problemphyxmini0228-displayedanswertorsionconstantsi, def:physics:phyx-mini-0228:phyxminiproblems-problemphyxmini0228-displayedanswertolerancesi}
  Agreement of a physical torsion constant with a displayed rounded choice
\end{definition}
\begin{proof}
  This is the declared model definition of Matches Answer Choice.
\end{proof}
\begin{lemma}[moment Of Inertia In SI eq one div two Million]
  \label{lem:physics:phyx-mini-0228:phyxminiproblems-problemphyxmini0228-momentofinertiainsi-eq-one-div-twomillion}
  \lean{PhyXMiniProblems.ProblemPhyXMini0228.momentOfInertiaInSI_eq_one_div_twoMillion}
  \uses{def:physics:phyx-mini-0228:phyxminiproblems-problemphyxmini0228-momentofinertiainsi, def:physics:phyx-mini-0228:phyxminiproblems-problemphyxmini0228-watchbalancewheelsetup, def:physics:phyx-mini-0228:phyxminiproblems-problemphyxmini0228-matchesproblemandfigurereadouts, def:physics:phyx-mini-0228:phyxminiproblems-problemphyxmini0228-satisfiesrimmomentofinertialaw}
  The source mass and radius, together with the rim-concentrated model, imply the intermediate moment of inertia I = 1 / 2,000,000 kg m\textasciicircum{}2
\end{lemma}
\begin{proof}
  The conclusion follows from the governing relations and figure readouts named in the statement of moment Of Inertia In SI eq one div two Million.
\end{proof}
% --- Archon declaration topology end ---
% --- Archon physics formalization source end ---
